\documentclass[lineno]{biometrika}

\usepackage{amsmath}
\usepackage{graphicx}
\usepackage{newtxtext}
\usepackage[subscriptcorrection]{newtxmath}
\usepackage{xr}
\externaldocument[S-]{supplement}

\def\Bka{{\it Biometrika}}
\newcommand{\pr}{\mathrm{pr}}
\newcommand{\LC}{\mathcal{L}}
\newcommand{\var}{\mathrm{var}}

\begin{document}

\jname{Biometrika}
\jyear{2026}
\jvol{}
\jnum{}
\cyear{2026}
\accessdate{}
%\received{}
%\revised{}
\makeatletter\def\printhistory{}\makeatother % no history line before acceptance

\markboth{K. W. Park}{Latent coverage from noisy calibration}

\title{Latent coverage from noisy calibration}

\author{KUN WOO PARK}
\affil{Department of Political Science and International Relations, Kookmin University,\\ 77 Jeongneung-ro, Seongbuk-gu, Seoul 02707, Republic of Korea \email{pkw31386094@gmail.com}}

\maketitle

\begin{abstract}
Prediction intervals for a latent quantity, such as the mean of an unsampled area or a response measured with error, are often calibrated on noisy proxies. We ask what latent coverage a threshold calibrated on noisy scores guarantees when the law and the variance of the noise are unknown. The answer turns on the shape of the latent law. If it is unimodal, even with mean zero, the guarantee is no better than with no shape assumption: with median-zero noise, noisy coverage $0.9$ guarantees only $0.8$. If it is bi-log-concave, which allows asymmetry and some multimodality, almost all of this loss is recovered for Gaussian, symmetric unimodal and mean-zero log-concave noise: under Gaussian noise, noisy coverage $0.9$ guarantees $0.8991$. When the noise laws and variances differ between units and are unknown, we give a conformal rank that keeps the latent guarantee, show that the usual high-probability rank can lose it as the calibration sample grows, and identify the smallest valid rank, exactly for symmetric unimodal noise and within a narrow range for Gaussian and log-concave noise. In simulations the guaranteed intervals were $16$--$48\%$ shorter than the rule for the same noise class without a latent shape assumption, and at most $3\%$ wider than the usual rank.
\end{abstract}

\begin{keywords}
Conformal prediction; Log-concavity; Measurement error; Small area estimation; Tolerance interval.
\end{keywords}

\section{Introduction}\label{sec:intro}

Suppose that a residual $W$, the latent target minus a centre fitted on independent training data, is observed only through $V = W + e$, with noise $e$ independent of $W$. A survey statistician who calibrates a prediction interval for the true mean of an unsampled area on the direct estimates of sampled areas is in this position \citep{fay1979}, as is a conformal analyst whose calibration responses carry measurement error \citep{einbinder2024, cohen2025}. The threshold $t$ at which the noisy scores $|V|$ have coverage $p$ is observable. We ask what latent coverage $\pr(|W| \le t)$ it guarantees, uniformly over a class of latent laws and over noise laws of every scale.

Two answers are known. If $W$ and $e$ are both symmetric and unimodal, the theorem of \citet{anderson1955} gives $\pr(|W + e| \le t) \le \pr(|W| \le t)$, so the noisy threshold is conservative \citep{einbinder2024}. Without a shape assumption, median-zero noise turns a noisy level $1 - \alpha$ into a latent level $1 - 2\alpha$ \citep[Proposition~2.3]{guille2024}; we show that this bound is attained, so it cannot be improved.

Related methods target different quantities. LatentCP \citep{zheng2026} gives a marginal latent guarantee through a known forward model, whereas we assume only a class of noise laws; \citet{cohen2025} deconvolve without a latent guarantee; conformal intervals for small areas cover observable responses \citep{bersson2024} or the expectations of sampled units \citep{zhang2026}. Empirical Bayes and parametric small-area intervals \citep{armstrong2022, ignatiadis2022, hallmaiti2006, yoshimori2014} concern sampled units or rely on a model for the latent effects.

We assume instead that the latent distribution is bi-log-concave \citep{dumbgen2017}, a class that contains the log-concave laws and some multimodal ones and does not require symmetry: residuals of area means of skewed outcomes can be asymmetric, and a mixture of two types of area can be bimodal and still bi-log-concave (Supplementary Lemma~\ref{S-lem:bimodal}); Anderson's theorem covers neither. Unimodality of the latent law, even with mean zero, gains nothing over no shape assumption (Proposition~\ref{prop:latentshape}); bi-log-concavity recovers almost all of the loss, at a price we give for each class of noise laws, exactly or within certified bounds (Theorem~\ref{thm:classes}, Table~\ref{tab:map}). The usual rank, the smallest whose noisy coverage holds with the prescribed probability, can lose the latent guarantee as the calibration sample grows (Proposition~\ref{prop:fail}); we identify the smallest rank that keeps it when the noise laws and variances differ between units and are unknown (\S~\ref{sec:procedure}). In simulations it cost at most $3\%$ of width over the usual rank and gave intervals $16$--$48\%$ shorter than the rule for the same noise class without a latent shape assumption (\S~\ref{sec:numerical}). For Gaussian noise of known variance, the location of the latent mean decides whether the required widening is of the order of the noise variance, as for a fixed smooth law \citep{jochmans2024}, or of its standard deviation (\S~\ref{sec:mean}).

\section{Coverage transfer}\label{sec:general}

\subsection{One end of the interval}

Let $\LC$ be the class of log-concave laws on the real line, including point masses, and $\mathcal B \supset \LC$ the class of laws that are point masses or have a continuous distribution function $F$ with $\log F$ and $\log(1 - F)$ concave. Both classes are closed under translation and scaling, and convolution with a log-concave law preserves them, by the theorem of \citet{prekopa1973} (Supplementary \S~\ref{S-sec:s-prelim}). Let $\mathcal E$ be a class of noise laws closed under scaling and negation, and let
\[
\Psi_{\mathcal E}(q) = \sup_{\epsilon \in \mathcal E} E\{\min(1, q \mathrm{e}^{-\epsilon})\}.
\]

The extremal latent laws are exponential tails, because concavity of $\log F_Y$ limits how fast mass can accumulate below a quantile: with $E_u$ exponential with rate $u > 0$, the law of $b - E_u$ has distribution function $\min\{1, \mathrm{e}^{u(y - b)}\}$.

\begin{proposition}\label{prop:general}
Let $\epsilon$ have a law in $\mathcal E$ and be independent of $Y$, with $\log F_Y$ concave. If $\pr(Y + \epsilon \le 0) \ge p > \Psi_{\mathcal E}(q)$, then $F_Y(0) \ge q$. If $p < \Psi_{\mathcal E}(q)$, some exponential tail $Y$ with $F_Y(0) < q$ and some $\epsilon$ with law in $\mathcal E$ have $\pr(Y + \epsilon \le 0) > p$.
\end{proposition}

\begin{proof}
Suppose that $v = F_Y^{-1}(q) > 0$, let $a = \log(1/q)$ and let $u$ be a supergradient of $\log F_Y$ at $v$. Then $F_Y(y) \le \min\{1, q\mathrm{e}^{u(y - v)}\}$, so the exponential tail $v + a/u - E_u$ lies stochastically below $Y$. Adding $\epsilon$ preserves this, and as $v > 0$ and $u\epsilon$ has a law in $\mathcal E$,
\[
p \le \pr(Y + \epsilon \le 0) \le \pr(a/u - E_u + \epsilon \le 0) = E\{\min(1, q \mathrm{e}^{-u\epsilon})\} \le \Psi_{\mathcal E}(q).
\] If $\log F_Y$ has no finite supergradient at $v$, as when $Y$ has an atom at $v$ and no mass below it, then $p \le \pr(\epsilon < 0) \le \lim_{u\to\infty} E\{\min(1, q \mathrm{e}^{-u\epsilon})\} \le \Psi_{\mathcal E}(q)$.

Conversely, $a - E_1$ with noise $\epsilon$ attains $E\{\min(1, q \mathrm{e}^{-\epsilon})\}$; shift it slightly to the right.
\end{proof}

Applied to $Y = (W - t)/s$ for any scale $s > 0$, Proposition~\ref{prop:general} controls the right end of $[-t, t]$; applied to $-W$, it controls the left end.

\subsection{Both ends}

A noisy failure probability $1 - p$ may be split in any way between the two ends, and each end transfers separately. Let $\lambda_{\mathcal E}(p) = \sup\{q' : \Psi_{\mathcal E}(q') < p\}$, the inverse of $\Psi_{\mathcal E}$ when $\Psi_{\mathcal E}$ is continuous and increasing, and $\phi_{\mathcal E}(a) = 1 - \lambda_{\mathcal E}(1 - a)$. With $\beta_{\mathcal E} = \inf_{\epsilon \in \mathcal E} \pr(\epsilon \ge 0)$, $\Psi_{\mathcal E}(q') \le 1 - (1 - q')\beta_{\mathcal E}$; we assume $\beta_{\mathcal E} > 0$, so that $\phi_{\mathcal E}(0) = 0$, and treat $\Psi_{\mathcal E} \equiv 1$, where no transfer is possible, separately (Supplementary \S~\ref{S-sec:s-oneend}). In words, $\phi_{\mathcal E}(a)$ is the largest latent failure probability at one end that is compatible with noisy failure probability $a$ there, by Proposition~\ref{prop:general}. With $a_L$ and $a_R$ the noisy failure probabilities at the left and right ends, $\pr(W + e < -t)$ and $\pr(W + e > t)$, define the sufficient noisy level
\begin{equation}\label{eq:pE}
p_{\mathcal E}(q) = \inf\{p : \phi_{\mathcal E}(a_L) + \phi_{\mathcal E}(a_R) \le 1 - q \text{ whenever } a_L + a_R \le 1 - p\}.
\end{equation}
Noisy coverage above $p_{\mathcal E}(q)$ at threshold $t$ then gives latent coverage $q$ for every $W \in \mathcal B$ and every noise with law in $\mathcal E$. Putting all the failure at one end, with small noise so that the other end is untouched, shows that no noisy level below $\Psi_{\mathcal E}(q)$ suffices (Supplementary \S~\ref{S-sec:s-oneend}). So $\Psi_{\mathcal E}(q)$ is the level needed when all the noisy failure sits at one end, and $\Psi_{\mathcal E}(q) \le p_{\mathcal E}(q)$. If $\phi_{\mathcal E}$ is convex, then, as $\phi_{\mathcal E}(0) = 0$, $\phi_{\mathcal E}(a_L) + \phi_{\mathcal E}(a_R) \le \phi_{\mathcal E}(a_L + a_R)$: the worst split puts the whole failure at one end, and $p_{\mathcal E}(q) = \Psi_{\mathcal E}(q)$. For example, with symmetric unimodal noise and noisy failure probability $0.1$, the latent failure probability can reach $\phi_{\mathcal E}(0.1) = 0.1019$ when all the noisy failure is at one end, but only $2\phi_{\mathcal E}(0.05) = 0.1009$ when it is split equally.

\begin{theorem}\label{thm:classes}
Let $q > \mathrm{e}^{-1}$.
\par\noindent (i) For symmetric unimodal noise, $\Psi_{\mathcal E}(q) = (1 + q \mathrm{e}^{-r})/2$, where $r > 0$ solves $(1 + r)\mathrm{e}^{-r} = (1 + \log q)/q$, attained by uniform noise; $\Psi_{\mathcal E}$ is strictly convex, so $\phi_{\mathcal E}$ is convex and $p_{\mathcal E}(q) = \Psi_{\mathcal E}(q)$.
\par\noindent (ii) For mean-zero log-concave noise, $\Psi_{\mathcal E}(q)$ is the supremum over centred log-affine laws on segments, and at $q = 0.9$ it lies in $(0.90517, 0.9052)$, the lower end from the centred exponential.
\par\noindent (iii) For mean-zero unimodal noise, $\Psi_{\mathcal E}(q) = 1$.
\end{theorem}

Part (i), where $q > \mathrm{e}^{-1}$ makes $1 + \log q > 0$, holds because symmetric unimodal laws are scale mixtures of uniform laws; the proof of convexity is analytic (Supplementary Proposition~\ref{S-prop:su-exact}). For part (ii), the localization theorem \citep{fradelizi2004} reduces the supremum to log-affine laws on segments; a branch and bound in ball arithmetic bounds it, and certified splits of the failure probability between the ends give $p_{\mathcal E}(0.9) \le 0.906$ and, at noisy level $0.9$, latent coverage $0.8935$. For part (iii), a noise spread uniformly over $[-a, 0]$ with $a$ large, so that $q\mathrm{e}^{-\epsilon} \ge 1$ except on $[\log q, 0]$, and balanced to mean zero by a small mass far above zero, makes $E\{\min(1, q\mathrm{e}^{-\epsilon})\}$ arbitrarily close to $1$. For Gaussian noise, $p_{\mathcal E}(0.9) \in (0.90079, 0.901]$ (\S~\ref{sec:mean}).

\begin{table}
\tbl{Latent coverage guaranteed by noisy coverage $0.9$, the noisy level that guarantees latent coverage $0.9$, and the smallest valid rank (Theorem~\ref{prop:unknown}) at $\delta = 0.05$}
{\begin{tabular}{llcccc}
latent law & noise & latent coverage & noisy level & $k$, $K = 10^3$ & $k$, $K = 10^4$ \\
symmetric unimodal & symmetric unimodal & $\ge 0.9$ & $0.9$ & 916 & 9050 \\
$\mathcal B$ & Gaussian & $0.8991$ & $(0.90079, 0.901]$ & 917 & 9058--9060 \\
$\mathcal B$ & symmetric unimodal & $0.8981$ & $0.9017873\ldots$ & 918 & 9068 \\
$\mathcal B$ & mean-zero log-concave & $0.8935$ & $(0.90517, 0.906]$ & 921--922 & 9101--9109 \\
$\mathcal B$ & mean-zero unimodal & none & $1$ & -- & -- \\
unimodal, mean zero & median zero & $0.8$ & $0.95$ & 962 & 9537 \\
unimodal, mean zero & mean-zero log-concave & $1 - \mathrm{e}/10$ & $1 - 0.1/\mathrm{e}$ & 974 & 9664 \\
\end{tabular}}
\label{tab:map}
\begin{tabnote}
Coverages and levels are certified. An interval of levels runs from a certified lower bound on $\Psi_{\mathcal E}(0.9)$, for Gaussian noise that of Corollary~\ref{cor:slack}, to a certified sufficient level; in a range of ranks the lower is necessary and the upper sufficient. The first row is valid by Anderson's theorem. Latent coverages are rounded down.
\end{tabnote}
\end{table}

\subsection{Unimodality does not suffice}

\begin{proposition}\label{prop:latentshape}
Let $\beta_{\mathcal E} > 0$ and $1 - \beta_{\mathcal E} < p < 1$. Then $\pr(|W| \le 1) \ge 1 - (1 - p)/\beta_{\mathcal E}$ for every $W$ and every noise with law in $\mathcal E$ and $\pr(|W + e| \le 1) \ge p$. If the infimum is approached by laws continuous at $0$, then for every $\gamma > 0$ there are a unimodal $W$ with $E(W) = 0$ and such a noise with $\pr(|W| \le 1) \le 1 - (1 - p)/\beta_{\mathcal E} + \gamma$. If $\beta_{\mathcal E} = 0$, approached in the same way, then for every $p < 1$ and $\gamma > 0$ there are such $W$ and noise with $\pr(|W| \le 1) \le \gamma$.
\end{proposition}

The lower bound is that of \citet[Proposition~2.3]{guille2024}, stated for a class of noise laws; what Proposition~\ref{prop:latentshape} adds is that it is attained by unimodal latent laws with mean zero. For unimodal latent laws the guarantee is therefore the shape-free one: $2p - 1$ for median-zero noise, $1 - \mathrm{e}(1 - p)$ for mean-zero log-concave noise, where $\beta_{\mathcal E} = \mathrm{e}^{-1}$ is attained by the centred exponential law, and none for mean-zero unimodal noise.

The extremal laws place a thin layer of mass just outside the threshold, and the noise carries a fraction $1 - \beta_{\mathcal E}$ of it inside. Bi-log-concavity excludes such a layer: concavity of $\log F$ makes $f/F$ nonincreasing, so the density cannot jump up at the threshold. Proposition~\ref{prop:latentshape} shows that unimodality cannot replace bi-log-concavity; it does not show that bi-log-concavity is necessary.


\section{Calibration with unknown noise}\label{sec:procedure}

Let $W_1, \ldots, W_K, W_{\rm new}$ be independent with law $G$, and let $e_i = \sigma_i\epsilon_i$, independent of them and of each other. The laws of the $\epsilon_i$ lie in $\mathcal E$ and may differ between units, and the scales $\sigma_i > 0$ are unknown. The scores $V_i = W_i + e_i$ are independent but not identically distributed. Let $\mathcal D$ denote the calibration data and $T = |V|_{(k)}$, the $k$th smallest of $|V_1|, \ldots, |V_K|$. We call $\pr_{\mathcal D}\{\pr(|W_{\rm new}| \le T \mid \mathcal D) \ge q\}$ the latent reliability of rank $k$: the probability, over calibration samples, that the interval attains latent coverage $q$.

\begin{theorem}\label{prop:unknown}
If $G \in \mathcal B$, then
\[
\pr_{\mathcal D}\{\pr(|W_{\rm new}| \le T \mid \mathcal D) \ge q\} \ge \pr\{{\rm Bin}(K, p_{\mathcal E}(q)) \le k - 1\}.
\]
Conversely, if $\pr\{{\rm Bin}(K, \Psi_{\mathcal E}(q)) \le k - 1\} < 1 - \delta$, there are a log-concave $G$ and a noise law in $\mathcal E$, the same in every unit, for which the left side is below $1 - \delta$. Hence the smallest rank valid for every $G \in \mathcal B$ and all such noise lies between the smallest $k$ with $\pr\{{\rm Bin}(K, \ell) \le k - 1\} \ge 1 - \delta$ at $\ell = \Psi_{\mathcal E}(q)$ and at $\ell = p_{\mathcal E}(q)$, which coincide when $p_{\mathcal E}(q) = \Psi_{\mathcal E}(q)$.
\end{theorem}

\begin{proof}
Let $Q_q$ denote the lower $q$-quantile and $r = Q_q(|W|)$ for $W \sim G$. With $N = \#\{i : |V_i| < r\}$,
\[
\pr(|W_{\rm new}| \le T \mid \mathcal D) \ge q \iff T \ge r \iff N \le k - 1.
\] For $s < r$, $\pr(|W| \le s) < q$, so by the transfer of \S~\ref{sec:general} at threshold $s$, $\pr(|V_i| \le s) \le p$ for every $p > p_{\mathcal E}(q)$, whatever the law and scale of $e_i$; letting $s \uparrow r$, $\pr(|V_i| < r) \le p_{\mathcal E}(q)$. $N$ is a sum of independent indicators with these success probabilities, so it is stochastically smaller than ${\rm Bin}\{K, p_{\mathcal E}(q)\}$.

Conversely, for $p < \Psi_{\mathcal E}(q)$ the one-ended laws of \S~\ref{sec:general} give a log-concave $G$ and a noise with $\pr(|V_i| < r) \ge \pr(|V_i| \le 1) > p$ and $r > 1$, so the left side is at most $\pr\{{\rm Bin}(K, p) \le k - 1\}$, which tends to $\pr\{{\rm Bin}(K, \Psi_{\mathcal E}(q)) \le k - 1\}$ as $p \uparrow \Psi_{\mathcal E}(q)$.
\end{proof}

The bound holds unit by unit, so the noise laws and variances need not be known or shared.

For a one-sided interval $(-\infty, t]$ the answer is exact for every noise class. By Proposition~\ref{prop:general}, noisy coverage $\pr(W + e \le t) > \Psi_{\mathcal E}(q)$ gives $\pr(W \le t) \ge q$ whenever $\log F_W$ is concave, and no lower noisy level suffices. The proof of Theorem~\ref{prop:unknown}, applied to the signed scores with $r = F_W^{-1}(q)$, then makes the smallest valid rank the smallest $k$ with $\pr\{{\rm Bin}(K, \Psi_{\mathcal E}(q)) \le k - 1\} \ge 1 - \delta$. At $q = 0.9$, $\Psi_{\mathcal E}(q)$ is $0.900803$ for Gaussian noise, $0.9017873\ldots$ for symmetric unimodal noise and lies in $(0.90517, 0.9052)$ for mean-zero log-concave noise.

To use Theorem~\ref{prop:unknown}, choose the noise class, take its certified sufficient level from Table~\ref{tab:map}, the upper end where the table gives an interval, for instance $0.901$ for Gaussian noise, and take as $k$ the smallest rank with $\pr\{{\rm Bin}(K, \text{level}) \le k - 1\} \ge 1 - \delta$. At $q = 0.9$ and $\delta = 0.05$ a valid rank exists for $K \ge 29$ under Gaussian or symmetric unimodal noise and $K \ge 31$ under mean-zero log-concave noise, against $K \ge 59$ and $K \ge 80$ without a latent shape assumption. For $K = 110$ the three bi-log-concave rows of Table~\ref{tab:map} give the usual rank, $k = 105$, against $109$ and $110$ without a shape assumption; the marginal rank $100$ is not valid. The raise grows with $K$: the required level exceeds $q$ by a fixed amount, while consecutive ranks differ in level by about $1/K$. It cannot be dropped:

\begin{proposition}\label{prop:fail}
Let $q \in (0,1)$. There are a log-concave law $G$ and $D > 0$ such that, with $e_i \sim N(0, D)$ for all $i$, $\pr_{\mathcal D}\{\pr(|W_{\rm new}| \le |V|_{(k_K)} \mid \mathcal D) \ge q\} \to 0$ as $K \to \infty$ for every sequence of ranks with $k_K/K \to q$, such as the marginal and, for fixed $\delta$, the usual high-probability rank.
\end{proposition}

\begin{proof}
By Theorem~\ref{thm:small} of \S~\ref{sec:mean}, whose constant $c_q$ is positive, some log-concave $W$ has $H(1) \ge q$ and $r_q = Q_q(|W|) > 1$, where $H$ is the distribution function of $|W + D^{1/2}Z|$, $Z \sim N(0, 1)$. As in the proof of Theorem~\ref{prop:unknown}, the latent reliability of rank $k$ is $\pr\{{\rm Bin}(K, H(r_q-)) \le k - 1\}$, and $H(r_q-) > H(1) \ge q$ as $H$ is strictly increasing, so it tends to zero by the law of large numbers.
\end{proof}

The latent coverage itself barely moves: at the extremal law of \S~\ref{sec:numerical} it tends to $0.89918$, a shortfall of $0.00082$. What fails is the guarantee: once $K$ is large enough for the order statistic to resolve this shortfall, the probability of latent coverage $q$ tends to zero; at $K = 10^6$ the usual rank has exact latent reliability $0.151$, against $0.989$ for the rank of Theorem~\ref{prop:unknown} with Gaussian noise. For $29 \le K \le 300$ the usual rank's reliability falls short of $0.95$ by less than $0.005$, so the raise matters most for large $K$.

\section{Gaussian noise}\label{sec:mean}

\subsection{The sharp radius}\label{sec:setup}

Gaussian noise is one of the classes of \S~\ref{sec:general}. When its variance $D$ is known, the noisy threshold can be widened instead of its rank raised, and the widening needed scales with $D$ in a way that depends on where the latent mean lies. After scaling $t$ to $1$, write $x = D/t^2$, let $Z \sim N(0,1)$ be independent of $W$ and define
\begin{equation}\label{eq:R}
R_{p,q}(x) = \sup\{ Q_q(|W|) : W \in \LC,\ \pr(|W + x^{1/2}Z| \le 1) \ge p \},
\end{equation}
with $R_{p,q}(x) = -\infty$ if nothing is feasible. If calibration certifies noisy coverage $p$ at threshold $t$, then $t\,R_{p,q}(D/t^2)$ is the smallest radius that guarantees latent coverage $q$ for every log-concave latent law. The radius is sharp for this class and this single constraint; we do not claim optimality among procedures that use the whole calibration sample. A reduction to log-affine laws on segments makes $R_{p,q}$ computable (Supplementary Proposition~\ref{S-prop:reduction}). Let $R^{\mathcal B}_{p,q} \ge R_{p,q}$ be defined with $\mathcal B$ in place of $\LC$.

The behaviour of $R_{p,q}$ for small $x$ is governed by a one-sided constant,
\[
c_{p,q} = \sup\{F_Y^{-1}(q) : Y \in \LC,\ \pr(Y + Z \le 0) \ge p\}, \qquad c_q = c_{q,q},
\]
which is attained along exponential tails and is the same over $\mathcal B$ (Supplementary Lemma~\ref{S-lem:onesided}). Ball arithmetic gives $c_{0.8} = 0.04607$, $c_{0.9} = 0.01906$ and $c_{0.95} = 0.00841$ to the digits shown.

\subsection{The widening and the location of the mean}

Without a restriction on the mean the widening is of order $x^{1/2}$.

\begin{theorem}\label{thm:small}
For every $q \in (0,1)$, every $x > 0$ and every $W \in \mathcal B$ with $\pr(|W + x^{1/2}Z| \le 1) \ge q$, $Q_q(|W|) \le 1 + c_q x^{1/2}$. Moreover $R_{q,q}(x) = 1 + c_q x^{1/2} + o(x^{1/2})$ as $x \to 0$, and likewise for $R^{\mathcal B}_{q,q}$.
\end{theorem}

\begin{proof}[of the upper bound]
Let $\alpha_R = \pr(W + x^{1/2}Z > 1)$ and $\alpha_L = \pr(W + x^{1/2}Z < -1)$, so that $\alpha_L + \alpha_R \le 1 - q$. The law of $Y = (W - 1)/x^{1/2}$ has $\log F_Y$ concave and $\pr(Y + Z \le 0) = 1 - \alpha_R$. By Supplementary Lemma~\ref{S-lem:onesided}, which holds over $\mathcal B$ and gives that $c_q$ is nonincreasing in $q$, $F_Y^{-1}(1 - \alpha_R) \le c_{1-\alpha_R} \le c_q$, that is, $\pr(W > 1 + c_q x^{1/2}) \le \alpha_R$. The same argument applied to $-W$, which uses concavity of $\log(1 - F_W)$, gives $\pr(W < -1 - c_q x^{1/2}) \le \alpha_L$, so $\pr(|W| \le 1 + c_q x^{1/2}) \ge q$.
\end{proof}

The lower bound uses exponential tails just outside the threshold; at $q = 0.9$ the widening is below $1.16\%$ of the radius for every $W \in \mathcal B$ (Supplementary \S~\ref{S-sec:s-feasible}).

A mean away from the boundary reduces the order to $D$: for a mean-zero latent law $Q_q(|W|) \le (t^2 + D)^{1/2}$ (Supplementary Theorem~\ref{S-thm:mean}, via the heat equation for quantiles \citep{steinbrecher2008}), an order that centred log-affine laws attain, so centring does not replace symmetry (Supplementary Proposition~\ref{S-prop:centred}). The order $D^{1/2}$ returns only when the mean is within a bounded number of noise standard deviations of the boundary (Supplementary Theorem~\ref{S-thm:transition}).

\subsection{Slack}\label{sec:slack}

Theorem~\ref{thm:small} gives a positive widening because $c_q > 0$. A noisy level slightly above $q$ removes the widening at every noise level. Splitting both the noisy and the latent failure probabilities between the two ends gives $Q_q(|W|) \le 1 + C_{p,q}x^{1/2}$ for every $x$ and every $W \in \mathcal B$ feasible at level $p \ge q$, with a constant $C_{p,q} \ge c_{p,q}$ defined and certified in the Supplementary Material.

\begin{corollary}\label{cor:slack}
Let $q = 0.9$. If $p \ge 0.901$, then $R^{\mathcal B}_{p,q}(x) \le 1$ for every $x > 0$; if $p \le 0.90079$, then $R_{p,q}(x) > 1$ for all small $x$. Hence for Gaussian noise $p_{\mathcal E}(0.9) \in (0.90079, 0.901]$, the Gaussian row of Table~\ref{tab:map}, with both ends certified in ball arithmetic. The same certificates with $p = 0.9$ fixed and $q$ varying (Supplementary \S~\ref{S-sec:s-ball}) show that noisy coverage $0.9$ gives latent coverage $0.8991$ for every $W \in \mathcal B$ and every noise variance, while $0.8992$ fails for some $W \in \LC$.
\end{corollary}

\label{sec:hetero}Known variances allow shorter intervals: in simulations a radius for log-concave latent laws, evaluated as a numerical upper bound, was $9.5$--$12\%$ shorter than the noisy threshold (Supplementary \S~\ref{S-sec:s-hetero}, Table~\ref{S-tab:synth}).

\section{Numerical illustration}\label{sec:numerical}

Table~\ref{tab:stress} shows the failure of Proposition~\ref{prop:fail} at the extremal law: an exponential tail whose end lies $2\%$ beyond the threshold, with Gaussian noise of equal variances and standard deviation $1\%$ of the threshold, located so that the population noisy coverage is $0.9$; its latent coverage is then $0.89918$. The usual rank keeps its noisy guarantee, with exact noisy reliability $0.952$ and $0.951$ at $K = 10^3$ and $10^4$, but its latent reliability falls below $0.95$, while the raised rank of Theorem~\ref{prop:unknown} keeps it.

\begin{table}
\tbl{Stress test at the extremal law: mean latent coverage given the data and latent reliability, target $0.9$ and $0.95$, $\delta = 0.05$}
{\begin{tabular}{llccc}
$K$ & rank & $k$ & mean latent coverage & latent reliability \\
$10^3$ & marginal & 901 & $0.89931$ & $0.482$ \\
 & usual & 916 & $0.91483$ & $0.943$ \\
 & raised, Gaussian noise & 917 & $0.91588$ & $0.954$ \\
$10^4$ & marginal & 9001 & $0.89918$ & $0.399$ \\
 & usual & 9050 & $0.90419$ & $0.918$ \\
 & raised, Gaussian noise & 9060 & $0.90522$ & $0.958$ \\
\end{tabular}}
\label{tab:stress}
\begin{tabnote}
Latent reliability exact, as in the proof of Proposition~\ref{prop:fail}; mean latent coverage over $20000$ ($K = 10^3$) and $10000$ ($K = 10^4$) data sets, standard errors at most $0.00007$. Further results in the Supplementary Material.
\end{tabnote}
\end{table}

Table~\ref{tab:sameinfo} shows what a latent guarantee costs, and what the latent shape assumption saves, when no rule is told the noise law within its class or the variances: four standardized log-concave latent laws (normal, Laplace, centred gamma with shape $2$, truncated exponential), five noise laws (Gaussian, uniform, Laplace, $t_3$, centred exponential) with heterogeneous variances, and exact latent coverage given each of $2000$ data sets per setting (Supplementary Material).

\begin{table}
\tbl{Rules with the same information: rank $k$ and mean width over the settings a rule covers, in percent above the usual rank and below the rule for the same noise class without a latent shape assumption; ranges, with medians in parentheses; target reliability $0.95$}
{\begin{tabular}{lcccccc}
 & \multicolumn{3}{c}{$K = 110$} & \multicolumn{3}{c}{$K = 1000$} \\
rule & $k$ & above usual & below no shape & $k$ & above usual & below no shape \\
usual & 105 & 0 & -- & 916 & 0 & -- \\
Gaussian & 105 & 0 & 22--29 (26) & 917 & 0.3--0.4 & 17--21 (21) \\
symmetric unimodal & 105 & 0 & 20--33 (27) & 918 & 0.6--1.0 & 16--24 (21) \\
mean-zero log-concave & 105 & 0 & 30--48 (39) & 922 & 2.0--2.9 & 21--33 (26) \\
no latent shape, median-zero noise & 109 & 25--49 & -- & 962 & 20--33 & -- \\
no latent shape, log-concave noise & 110 & 42--94 & -- & 974 & 28--54 & -- \\
\end{tabular}}
\label{tab:sameinfo}
\begin{tabnote}
Noisy levels $0.9$, $0.901$, $\Psi_{\mathcal E}(0.9)$, $0.906$, $0.95$ and $1 - 0.1/\mathrm{e}$. A rule covers a setting when the noise is in its class: Gaussian; symmetric unimodal (Gaussian, uniform, Laplace, $t_3$); mean-zero log-concave (Gaussian, uniform, Laplace, centred exponential); the first three need a bi-log-concave latent law. The Gaussian and symmetric unimodal rules are compared with the median-zero rule, the log-concave rule with the log-concave one. The usual rank has no guarantee uniform over bi-log-concave latent laws; Anderson's theorem covers it only for symmetric unimodal laws.
\end{tabnote}
\end{table}

Every rule had reliability at least $0.998$ where its class applies, and so did the usual rank in all $20$ settings: these laws are far from the extremal one of Table~\ref{tab:stress}, so they measure the cost of the guarantee, not its need. The guarantee is therefore cheap, and the shape assumption is what makes it cheap: without it the noisy levels $0.95$ and $1 - 0.1/\mathrm{e}$ cannot be lowered, by Proposition~\ref{prop:latentshape}.

The rule of Theorem~\ref{prop:unknown} needs no noise variances. In the school-district application of the Supplementary Material, with $K = 110$ and poorly estimated design variances, its interval covered at least $90\%$ of the remaining districts in all $480$ replications and was $19$--$26\%$ wider than a Fay--Herriot tolerance interval that plugs in the variances and fell short in up to $9\%$ of replications; the noise class and the shared latent law cannot be confirmed from these data.

\section{Discussion}\label{sec:discussion}

The guarantees hold with probability $1 - \delta$ over calibration samples: on that event, the latent coverage of the new unit, given the calibration data, is at least $q$. They allow heterogeneous noise within the specified classes, independent of the latent residual, but require a bi-log-concave latent law shared by the calibration units and the new one, the harder assumption to check. If the new residual is shifted by at most $\rho$ latent standard deviations and the latent law is log-concave, the latent coverage can fall by at most $\rho$, and by $q(1 - \mathrm{e}^{-\rho}) = q\rho + O(\rho^2)$ in the limit for some laws, so the linear order cannot be improved (Supplementary Proposition~\ref{S-prop:shift}). By comparison, at $q = 0.9$ Gaussian, symmetric unimodal and mean-zero log-concave noise cost at most $0.0065$ in coverage.

Bi-log-concavity cannot be checked on the latent scale, but convolution with log-concave noise preserves it: with a common log-concave noise law, for instance Gaussian noise with equal variances, a confidence band for the signed noisy residuals that contains no bi-log-concave distribution function \citep{dumbgen2017} rules it out, though not conversely.

Open questions include convexity of $\phi_{\mathcal E}$ for Gaussian and log-concave noise, which would make their two-sided levels equal to $\Psi_{\mathcal E}(q)$ as in Theorem~\ref{thm:classes}(i) and narrow the ranges of ranks in Table~\ref{tab:map}, at most $8$ ranks wide at $K = 10^4$; an analytic proof that the centred exponential law is extremal among mean-zero log-concave noise laws; and whether $R^{\mathcal B}_{p,q} = R_{p,q}$.

\section*{Declaration of the use of generative AI and AI-assisted technologies}

The author used Claude (Anthropic) only to check grammar and English usage in the final draft. All results, proofs, computations and text are the author's own, and the author takes full responsibility for the content.

\section*{Supplementary material}\label{SM}

The Supplementary Material contains the omitted proofs, the certificates, further simulations and an application. Code reproducing every number, at tag \texttt{biometrika-submission}, is at\\ \texttt{https://github.com/kunhailP/latent-coverage-noisy-thresholds}.

\bibliographystyle{biometrika}
\bibliography{refs}

\end{document}
